% Source-aligned English for printed p.182.
\SourcePage{182}

\LayerTarget{tohoku:proof:proposition-3.10.1:p2}
% The incoming continuation is already integrated at the p.181 reading seam.

\LayerTarget{tohoku:corollary:3.10.1:1}
\noindent\textsc{Corollary 1.} \textit{Let $X$ be a space equipped with a
paracompactifying family $\Phi$. Then, for every abelian sheaf $F$ on $X$,}

\LayerTarget{tohoku:eq:corollary-3.10.1-1:support-limit}
\[
H^p_{\Phi}(X,F)=\varinjlim_U H^p(X,F_U),
\]
\textit{where the inductive limit is taken over the filtered ordered set of
open subsets $U$ of $X$ whose closure belongs to $\Phi$. (Here $F_U$ denotes
the sheaf on $X$ whose restriction to $U$ is $F|U$ and whose restriction to
$C U$ is zero.)}

\LayerTarget{tohoku:proof:corollary-3.10.1-1}
By Theorem 3.5.1,
$H^p(X,F_U)=H^p_{\Phi_U}(U,F)$, where $\Phi_U$ is the family of subsets of
$U$ that are closed in $X$. Setting
$\Gamma_{\Phi_U}(F)=H^0(X,F_U)$, we may also write
$H^p(X,F_U)=\mathbf R^p\Gamma_{\Phi_U}(F)$, in view of Proposition 3.1.3.
Hence Proposition 3.10.1 gives
$\varinjlim_U H^p(X,F_U)=\mathbf R^p\Gamma_{\Phi}(F)$, since
$\varinjlim_U\Gamma_{\Phi_U}(F)=\Gamma_{\Phi}(F)$, q.e.d.

\LayerTarget{tohoku:remark:corollary-3.10.1-1:cartan}
The preceding corollary is sometimes useful for reducing cohomology ``with
supports in $\Phi$'' to cohomology with arbitrary supports; it was pointed
out to me by H. Cartan.

\LayerTarget{tohoku:corollary:3.10.1:2}
\noindent\textsc{Corollary 2.} \textit{Let $X$ be a topological space and let
$Y$ be a subset of $X$ having a fundamental system of paracompact
neighborhoods. (It is enough, for example, that $X$ be metrizable, or locally
compact and paracompact.) Then, for every abelian sheaf $F$ on $X$,}

\LayerTarget{tohoku:eq:corollary-3.10.1-2:neighborhood-limit}
\[
H^p(Y,F)=\varinjlim_U H^p(U,F),
\]
\textit{where the limit is taken over the decreasing filtered ordered set of
open neighborhoods $U$ of $Y$ in $X$.}

\LayerTarget{tohoku:proof:corollary-3.10.1-2}
Indeed, the hypothesis gives
$H^0(Y,F)=\varinjlim_U H^0(U,F)$ ([11], Proposition 2.2.1), while the derived
functors of $F\longmapsto H^0(U,F)$ are the $H^p(U,F)$. Thus Corollary 2 is a
special case of the proposition. Observe also that
$H^0(Y,F)=\varinjlim_U H^0(U,F)$, and consequently Corollary 2 remains valid
when $Y$ is \textit{closed} and has \textit{a paracompact neighborhood}; the
proof is analogous to that of [11]. The cited work also gives a simple
counterexample, with $p=0$, for the case in which no paracompactness
hypothesis is imposed.

\LayerTarget{tohoku:result:3.10.2:preliminary}
As a supplement, we record without proof the following result, a special case
of general results on projective systems. Let $X$ be a locally compact space,
and consider the increasing filtered set of relatively compact open subsets
$U$ of $X$. For every abelian sheaf $F$ on $X$, the restriction homomorphisms
$H^p(X,F)\longrightarrow H^p(U,F)$ define canonical homomorphisms, which are
evidently homomorphisms of $\partial$-functors:

\LayerTarget{tohoku:eq:3.10.2}
\[
\tag{3.10.2}
H^p(X,F)\longrightarrow\varprojlim_U H^p(U,F),
\]
which are evidently bijective for $p=0$.

\TohokuContinuation{183}
